% Secondary table, shared by every experiments organization. A compact
% per-condition view of the primary arm against the strongest baseline, so a
% reader can see whether the aggregate hides a reversal. Booktabs only; the
% Difference row is meaningful only because the pairing is exact.
% One column per condition. With short condition keys (split seeds) this fits a
% column; with benchmark names it does not, and an overfull \hbox in a CVPR
% column is what you get. acwidetable is a plain table in a single-column
% layout and table* in cvpr2, so this costs nothing where it already fitted.
\begin{acwidetable}[t!]
  % Write your own caption. The lead phrase named this table in all five of the
  % first batch's papers because it used to be hardcoded here; a caption that
  % every paper in a proceedings shares is a fingerprint, not a convention.
  \caption{\textbf{\slot{name this comparison in your own words} (\slot{\up\ or \down}).}
    \slot{what each column is, which direction favors which arm, and one
    sentence on the estimator and repetition count}.}
  \label{tab:paired}
  \vspace{-0.2cm}
  \centering
  {\fontsize{8pt}{10pt}\selectfont
   \begin{tabular}{l c c c c}
     \toprule
     & \textbf{\slot{c1}} & \textbf{\slot{c2}} & \textbf{\slot{c3}} & \textbf{\slot{c4}} \\
     \midrule
     \armbaseline & \slot{v} & \slot{v} & \slot{v} & \slot{v} \\
     \armprimary  & \slot{v} & \slot{v} & \slot{v} & \slot{v} \\
     \midrule
     Difference   & \slot{d} & \slot{d} & \slot{d} & \slot{d} \\
     \bottomrule
   \end{tabular}}
\end{acwidetable}
